% An invented paper for the end-to-end example: the study, its numbers and its reference do not exist.
\documentclass{article}
\begin{document}
We record an array of 24 thermometers once a minute, 1440 times a day, for 30 days.
Each thermometer therefore yields 43200 readings.
Across the array this gives 1036800 readings.
After removing 5\% of the readings as outliers, 984460 readings remain.
We model the temperature reported by a drifting thermometer as
\begin{equation}
  T(t) = T_0 + \alpha t \label{eq:model}
\end{equation}
so that over an interval $\tau$ the reported temperature changes by
\begin{equation}
  T(t + \tau) - T(t) = \alpha \tau \label{eq:drift}
\end{equation}
The mean fitted drift rate is 0.02 K per day, so the mean drift over the 30 days is 0.6 kelvin.
Lee reports that drifts of this size are typical of uncalibrated thermistors~\cite{lee2019}.
These results suggest that the array needs recalibration every month.
\begin{thebibliography}{1}
\bibitem{lee2019} C. Lee, Thermistor ageing in field arrays, J. Instr. Field Methods 7 (2019) 12.
\end{thebibliography}
\end{document}
